\exercisesection{3}

¶ 1. (a) Let L be a field and T an indeterminate; in the field of formal power series L((T)), consider the series

\[
s = c _ {0} + c _ {1} \mathrm{T} ^ {1!} + c _ {2} \mathrm{T} ^ {2!} + \dots + c _ {n} \mathrm{T} ^ {n!} + \dots
\]

where the \(c_{i} \in L\) are all \(\neq 0\) . Show that s is transcendental over the subfield \(\mathrm{L}(\mathrm{T})\) of \(\mathrm{L}((\mathrm{T}))\) . (Argue by reductio ad absurdum by supposing that there exists a relation \(a_{0}(\mathrm{T}) + a_{1}(\mathrm{T})s + \cdots + a_{q}(\mathrm{T})s^{q} = 0\) where \(a_{i}(\mathrm{T})\) are polynomials in \(\mathrm{L}(\mathrm{T})\) such that \(a_{q} \neq 0\) ; construct the equation satisfied by

\[
s ^ {\prime} = s - c _ {0} - c _ {1} \mathrm{T} ^ {1!} - \dots - c _ {n - 1} \mathrm{T} ^ {(n - 1)!}
\]

and show that for \(n\) sufficiently large its constant term is a polynomial of degree \(< n!\) , whence our contradiction.)

\begin{enumerate}
\def\labelenumi{(\alph{enumi})}
\setcounter{enumi}{1}
\tightlist
\item
  Let \(k\) be a field, \(\mathbf{X}, \mathbf{Y}, \mathbf{Z}\) , \(\mathrm{T}\) four indeterminates, \(\mathbf{K} = k(\mathbf{X}, \mathbf{Y}, \mathbf{Z})\) and \(\mathbf{E}\) an algebraic closure of \(k(\mathbf{X})\) ; in \(\mathbf{E}\) , for all \(n\) , let \(x_{n}\) denote an element such that \(x_{n}^{n} = \mathbf{X}\) . In the field of formal power series \(\mathrm{E}((\mathrm{T}))\) , consider the element
\end{enumerate}

\[
s = x _ {1} \mathrm{T} + x _ {2} \mathrm{T} ^ {2!} + \dots + x _ {n} \mathrm{T} ^ {n!} + \dots
\]

The elements \(s\) and \(\mathbf{T}\) are algebraically independent over \(\mathbf{E}\) by (a). Consider the homomorphism \(f\colon \mathbf{K} \to \mathbf{E}((\mathbf{T}))\) such that \(f(\mathbf{X}) = \mathbf{X}, f(\mathbf{Y}) = \mathbf{T}, f(\mathbf{Z}) = s\) . If \(v\) is the discrete valuation on \(\mathbf{E}((\mathbf{T}))\) defined by the order of formal power series (no. 3, Example 3), \(v \circ f\) is a discrete valuation on \(\mathbf{K}\) ; show that the residue field of \(v \circ f\) is the subfield \(k(x_1, x_2, \ldots, x_n, \ldots)\) of \(\mathbf{E}\) and is therefore infinite over \(k\) .

\begin{enumerate}
\def\labelenumi{\arabic{enumi}.}
\setcounter{enumi}{1}
\tightlist
\item
  Let \(\Gamma\) be a totally ordered group and \(k\) a field.
\end{enumerate}

\begin{enumerate}
\def\labelenumi{(\alph{enumi})}
\item
  Let A, B be two subsets of \(\Gamma\) well ordered by the induced ordering; show that the set \(A + B\) is well ordered and that, for all \(\gamma \in A + B\) , the set of ordered pairs \((\alpha, \beta) \in A \times B\) such that \(\alpha + \beta = \gamma\) is finite. (To prove the first assertion, argue by reductio ad absurdum by showing that otherwise it would be possible to define a strictly decreasing sequence \((\gamma_n)\) of elements of \(A + B\) such that \(\gamma_n = \alpha_n + \beta_n\) , where \(\alpha_n \in A\) , \(\beta_n \in B\) , the sequence \((\alpha_n)\) is strictly increasing and the sequence \((\beta_n)\) is strictly decreasing; conclude using Set Theory, Chapter III, §4, Exercise 3.)
\item
  Let \(S(\Gamma, k)\) be the set of families \(x = (x_{\alpha})_{\alpha \in \Gamma}\) such that \(x_{\alpha} \in k\) for all \(\alpha \in \Gamma\) and the set of \(\alpha\) for which \(x_{\alpha} \neq 0\) is a well ordered subset of \(\Gamma\) . Show that \(S(\Gamma, k)\) is an additive subgroup of \(k^{\Gamma}\) ; moreover, if \(x = (x_{\alpha}), y = (y_{\alpha})\) are two elements of \(S(\Gamma, k)\) , the set \(C\) of \(\alpha + \beta\) for the ordered pairs \((\alpha, \beta)\) such that \(x_{\alpha} \neq 0\) and \(y_{\beta} \neq 0\) is well ordered by (a) and, for all \(\gamma \in C\) , the element \(z_{\gamma} = \sum_{\alpha + \beta = \gamma} x_{\alpha} y_{\beta}\) is defined in \(k\) ; if we write \(z_{\alpha} = 0\) for \(\alpha \notin C\) , the element \(z = (z_{\alpha}) \in k^{\Gamma}\) belongs to \(S(\Gamma, k)\) ; it is denoted by \(xy\) . Show that with this law of composition and addition on the product group \(k^{\Gamma}\) the set \(S(\Gamma, k)\) is a field; moreover, for all \(x = (x_{\alpha}) \neq 0\) of \(S(\Gamma, k)\) , let \(\lambda\) be the least of the \(\alpha \in \Gamma\) such that \(x_{\alpha} \neq 0\) ; if we write \(v(x) = \lambda\) (and \(v(0) = +\infty\) ), show that \(v\) is a valuation on \(S(\Gamma, k)\) with \(k\) as residue field and \(\Gamma\) as value group. \(v\) is called the canonical valuation on the field \(S(\Gamma, k)\) . The field \(k\) is canonically identified with a subfield of \(S(\Gamma, k)\) . For all \(\alpha \in \Gamma\) , let \(u_{\alpha}\) denote the element \((x_{\lambda})_{\lambda \in \Gamma}\) such that \(x_{\alpha} = 1, x_{\lambda} = 0\) for \(\lambda \neq \alpha\) ; for all non-zero \(z \in S(\Gamma, k)\) , there is then a unique ordered pair \((t, \alpha) \in k \times \Gamma\) such that \(v(z) = \alpha\) and \(v(z - tu_{\alpha}) > \alpha\) . \(tu_{\alpha}\) is called the dominant term of \(z\) .
\end{enumerate}

¶ 3. Let K be a field, w a valuation on K, Γ the value group of w and k the residue field of w. We form the corresponding field S(Γ, k) and preserve the notation of Exercise 2. For all \(t \in k\), let \(t^{0}\) be an element of the class t in the valuation ring of w; for all \(\alpha \in \Gamma\), let \(u_{\alpha}^{0}\) be an element of K such that \(w(u_{\alpha}^{0}) = \alpha\).

\begin{enumerate}
\def\labelenumi{(\alph{enumi})}
\tightlist
\item
  Let M be a subset of S(Γ, k) containing the elements \(tu_{\alpha}\) for every ordered pair \((t, \alpha) \in k \times \Gamma\) ; suppose that there exists a bijection \(x \mapsto x^{0}\) of M onto a subset \(M^{0}\) of K and that the following conditions are fulfilled: (1) \((tu_{\alpha})^{0} = t^{0}u_{\alpha}^{0}\) for every ordered pair \((t, \alpha) \in k \times \Gamma\) ; (2) if \(x = (x_{\alpha})\) belongs to M and \(C_{x}\) is the well-ordered subset of Γ consisting of the α such that \(x_{\alpha} \neq 0\) , then for every segment D of \(C_{x}\) the element \((x_{\lambda})_{\lambda \in D}\) belongs to M; (3) if x, y are two elements of M and \(tu_{\alpha}\) is the dominant term of x - y, then \(w(x^{0} - y^{0} - t^{0}u_{\alpha}^{0}) > w(x^{0} - y^{0})\) . If \(z'\) is an element of K not belonging to \(M^{0}\) , show that there exists an element \(z \notin M\) in S(Γ, k) such that the mapping which coincides with \(x \mapsto x^{0}\) on M and maps \(z'\) to z also satisfies the above conditions. (Note that, if \(w(z') = \alpha\) , there exists a \(t^{0}\) such that
\end{enumerate}

\[
w (z ^ {\prime} - t ^ {0} u _ {\alpha} ^ {0}) > \alpha ;
\]

write \(z_{\alpha} = tu_{\alpha}\) , then show by transfinite induction that the \(z_{\beta}\) of index \(\beta > \alpha\) may be determined so that \(z = (z_{\lambda})_{\lambda \in \Gamma}\) solves the problem.)

\begin{enumerate}
\def\labelenumi{(\alph{enumi})}
\setcounter{enumi}{1}
\tightlist
\item
  Deduce from (a) that \(\operatorname{Card}(\mathbf{K}) \leqslant \operatorname{Card} S(\Gamma, k)\) .
\end{enumerate}

\begin{enumerate}
\def\labelenumi{\arabic{enumi}.}
\setcounter{enumi}{3}
\tightlist
\item
  Let A be a local integral domain, K its field of fractions, m its maximal ideal and v a valuation on K whose ring B dominates A. Suppose that, either m is a finitely generated ideal, or v is a discrete valuation; then there exists in m an element x such that \(v(x) = \inf_{y \in \mathfrak{m}} v(y)\) . Let \(A'\) be the subring of K generated by
\end{enumerate}

A and the elements \(yx^{-1}\) , where y runs through m; let \(A_{1}\) be the local ring of \(A'\) relative to the prime ideal \(p \cap A'\) , where p is the ideal of the valuation v. Show that the ring \(A_{1}\) does not depend on the choice of the element x satisfying the above conditions ( \(A_{1}\) is called ``the prime quadratic transform of A in the direction of \(v'\) ); if A is Noetherian, so is \(A_{1}\) .

¶ 5. Let k be a field and K = k(X, Y) the field of rational functions in two indeterminates over k. We write \(x_{1} = X\) , \(y_{1} = Y\) and for \(n \geqslant 1\) define inductively the elements \(x_{n+1} = y_{n}\) and \(y_{n+1} = x_{n}y_{n}^{-1}\) of K.

\begin{enumerate}
\def\labelenumi{(\alph{enumi})}
\tightlist
\item
  Write \(\mathbf{A}_n = k[x_n, y_n]\) ; the sequence \((\mathbf{A}_n)\) is increasing and, if
\end{enumerate}

\[
\mathfrak {m} _ {n} = \mathrm{A} _ {n} x _ {n} + \mathrm{A} _ {n} y _ {n},
\]

\(m_{n}\) is a maximal ideal of \(A_{n}\) and \(m_{n} = A_{n} \cap m_{n+1}\) ; in the ring \(A = \bigcap_{n} A_{n}\) , \(m = \bigcap_{n} m_{n}\) is a maximal ideal.

\begin{enumerate}
\def\labelenumi{(\alph{enumi})}
\setcounter{enumi}{1}
\item
  Let \(v\) be the valuation on \(\mathbf{K}\) defined by \(v(\mathbf{X}) = \rho = (1 + \sqrt{5}) / 2\) , \(v(\mathbf{Y}) = 1\) ; let \(\mathbf{B}\) and \(\mathfrak{p}\) be the ring and ideal of \(v\) ; show that \(\mathfrak{m} = \mathfrak{p} \cap \mathbf{A}\) (note that \(\rho^2 - \rho - 1 = 0\) and use this fact to calculate \(v(y_n)\) ). Show that for a monomial \(\mathbf{X}^i\mathbf{Y}^j\) ( \(i, j\) positive or negative integers) to belong to \(\mathbf{A}\) , it is necessary and sufficient that \(i\left(\frac{1 + \sqrt{5}}{2}\right) + j > 0\) ; for every monomial \(z = \mathbf{X}^i\mathbf{Y}^j\) , therefore, either \(z \in \mathbf{A}\) or \(1/z \in \mathbf{A}\) .
\item
  Show that \(B = A_{m}\) (write an element \(t \in B\) in the form of a quotient of two polynomials in X and Y and show in each of them the monomials at which v takes the least value).
\item
  Show that, for all \(n\) , \((\mathbf{A}_{n+1})_{\mathfrak{m}_{n+1}}\) is the prime quadratic transform of \((\mathbf{A}_n)_{\mathfrak{m}_n}\) in the direction of \(v\) (Exercise 4).
\end{enumerate}

\begin{enumerate}
\def\labelenumi{\arabic{enumi}.}
\setcounter{enumi}{5}
\tightlist
\item
  \begin{enumerate}
  \def\labelenumii{(\alph{enumii})}
  \tightlist
  \item
    Let \(\mathbf{A}\) be a valuation ring. Extend to finitely generated A-modules the results of Algebra, Chapter VII, § 4, nos. 1 to 4 (theory of invariant factors).
  \end{enumerate}
\end{enumerate}

\begin{enumerate}
\def\labelenumi{(\alph{enumi})}
\setcounter{enumi}{1}
\item
  Let K be a field, v a valuation on K, A the ring of v and \(\Gamma\) the order group of v. Let \(U = (\alpha_{ij})\) be a square matrix of order n with elements in A and of determinant 0; show that there exists \(\lambda \in \Gamma\) such that, for every square matrix \(U' = (\alpha_{ij}')\) of order n with elements in A and satisfying the conditions \(v(\alpha_{ij}' - \alpha_{ij}) > \lambda\) for every ordered pair \((i,j)\) , the invariant factors of \(U'\) are the same as those of U.
\item
  Under the hypotheses of (b), generalize the following results of Algebra, Chapter IX, § 5, no. 1, Theorem and Exercise 1 and § 6, Exercise 10 (for the last, suppose that K is not of characteristic 2 and that \(v(2) = 0\) ).
\end{enumerate}

\begin{enumerate}
\def\labelenumi{\arabic{enumi}.}
\setcounter{enumi}{6}
\item
  Show that the integral domain B defined in Chapter II, § 3, Exercise 2 (b) is a local ring whose maximal ideal is principal but is not a valuation ring.
\item
  Show that, if a valuation ring A is strongly Laskerian (Chapter IV, § 2, Exercise 28), A is a field or a discrete valuation ring (use Exercise 29 of Chapter IV, § 2).
\end{enumerate}
% semantic-fragments: legacy-input-seam
\relax{}
